\documentclass[UTF8,a4paper,zihao=-4,fontset=windows]{ctexart}

\usepackage{geometry}
\usepackage{setspace}
\usepackage{booktabs}
\usepackage{longtable}
\usepackage{tabularx}
\usepackage{array}
\usepackage{enumitem}
\usepackage{amsmath}
\usepackage{amssymb}
\usepackage{xcolor}
\usepackage{hyperref}

\geometry{left=2.8cm,right=2.6cm,top=2.6cm,bottom=2.6cm}
\setstretch{1.45}
\setlength{\parindent}{2em}
\setlength{\parskip}{0.25em}
\setlength{\emergencystretch}{3em}
\sloppy
\setlist{nosep,leftmargin=2.2em}
\hypersetup{
  colorlinks=true,
  linkcolor=black,
  urlcolor=blue,
  citecolor=black,
  pdfauthor={姜浩然},
  pdftitle={面向长时程多智能体大语言模型系统的上下文治理与长期记忆图谱方法研究开题报告}
}

\ctexset{
  section = {
    format = \zihao{3}\heiti\raggedright,
    name = {,、},
    number = \chinese{section},
    aftername = \hspace{0.2em}
  },
  subsection = {
    format = \zihao{4}\heiti\raggedright
  },
  subsubsection = {
    format = \zihao{-4}\heiti\raggedright
  }
}

\newcommand{\reporttitle}{面向长时程多智能体大语言模型系统的上下文治理与长期记忆图谱方法研究}
\newcommand{\framework}{象群-星鸦框架}
\newcommand{\blankline}{\underline{\hspace{5.2cm}}}

\begin{document}

\begin{titlepage}
  \centering
  \vspace*{1.2cm}
  {\zihao{1}\heiti 硕士学位论文开题报告\par}
  \vspace{1.1cm}
  {\zihao{2}\heiti 面向长时程多智能体大语言模型系统的\par}
  \vspace{0.25cm}
  {\zihao{2}\heiti 上下文治理与长期记忆图谱方法研究\par}
  \vspace{2.0cm}
  \begin{tabular}{>{\heiti}p{3.2cm}p{8cm}}
    学生姓名： & 姜浩然\\[0.45cm]
    学号： & \blankline\\[0.45cm]
    专业： & \blankline\\[0.45cm]
    研究方向： & 大语言模型智能体、上下文工程与多智能体系统\\[0.45cm]
    指导教师： & \blankline\\[0.45cm]
    所在学院： & \blankline\\[0.45cm]
  \end{tabular}
  \vfill
  {\zihao{4}二〇二六年四月\par}
\end{titlepage}

\begin{abstract}
随着大语言模型智能体从单轮问答逐步发展为能够规划、检索、调用工具、跨角色协作并持续保持状态的长时程系统，智能体系统的关键问题已经不再只是单次回答质量，而是如何在有限上下文预算、跨轮任务状态、跨智能体协作和长期经验复用之间保持稳定治理。现有自主智能体、多智能体系统、长上下文建模、提示压缩、检索增强生成和长期记忆研究为该问题提供了重要基础，但仍较少将运行内上下文治理、跨运行长期记忆组织、预算约束下的提示构造与事件中心可观测性统一建模。

本课题拟基于现有研究稿件《\reporttitle》展开硕士论文研究，将原统一稿件转化为一个总论与两个机制章节相互支撑的论文结构。研究拟提出\framework：其中，象群上下文管理策略负责运行内多智能体协作中的显著性排序、证据保留、上下文压缩和跨角色传递；克拉克星鸦长时记忆图谱机制负责将高价值执行经验提升为长期记忆胶囊，组织为包含任务、证据、时间、冲突和修正关系的记忆图谱，并在未来任务中以局部子图或摘要形式安全返回；MGCM 记忆治理层负责候选记忆的提升、验证、降级和废弃；长时程注意力、预算感知压缩和事件中心可观测性共同支撑系统连续执行、提示构造和失败诊断。

论文计划围绕五类研究问题展开实验：统一框架在长时程任务上的整体收益，象群上下文治理在高干扰上下文中的证据保留能力，长期记忆图谱在 recall、reuse、陈旧记忆抑制和冲突消解方面的表现，长时程注意力与上下文压缩在预算收缩下的鲁棒性，以及事件中心可观测性对任务轨迹重建和故障定位的作用。预期形成一套面向长时程多智能体大语言模型系统的上下文治理与长期记忆方法框架，为智能体底层编排框架、长时程任务执行和可解释系统评测提供可复现的理论建模、工程实现和实验依据。

\textbf{关键词：} 大语言模型智能体；多智能体系统；上下文治理；长期记忆图谱；象群上下文管理策略；克拉克星鸦长时记忆；MGCM；事件中心可观测性
\end{abstract}

\tableofcontents
\newpage

\section{本课题的研究目的与意义}

\subsection{研究目的}

本课题从当前研究稿件 \texttt{cas-sc-template-zh.tex} 出发，拟将“面向长时程多智能体大语言模型系统的上下文治理与长期记忆图谱方法”整理为硕士论文研究。论文不把智能体框架简单视为工具链堆叠，而是从第一性原理出发，将长时程多智能体系统抽象为在有限上下文预算下对信息、记忆、协作和事件证据进行持续治理的问题。研究目的主要包括以下几个方面：

\begin{enumerate}[label=\arabic*.]
  \item 建立长时程多智能体大语言模型系统的问题表述。将用户请求、线程历史、工具集合、检索结果、局部记忆、候选共享记忆、集体记忆和长期图谱记忆统一纳入候选上下文集合，明确提示构造本质上是受预算约束的上下文选择问题。
  \item 提出象群启发的多智能体上下文治理方法。针对多智能体协作中工具观测、委派数据包、中间产物和历史约束共同争夺提示空间的问题，设计一种以验证程度、复用价值、结果关联和冗余抑制为核心的系统级显著性策略。
  \item 提出克拉克星鸦启发的长期记忆图谱机制。将跨运行经验表示为可治理的记忆胶囊节点，通过时间相邻、同任务、证据支持、范围重叠、修正替代和显式冲突等关系构建长期记忆图谱，使长期记忆不再停留在平面文本累积或普通向量召回层面。
  \item 构建长时程运行的系统融合机制。将象群上下文治理、长期记忆图谱、MGCM 记忆治理、长时程注意力、预算感知上下文压缩和事件中心可观测性组织为统一框架，形成可插拔、可评估、可复现的智能体运行时方法。
  \item 通过公开 benchmark 切片和模块级强基线实验验证方法有效性。围绕 GAIA、HotpotQA、MuSiQue、LoCoMo、LongMemEval、LongBench、InfiniteBench 以及 AgentOrch 诊断套件设计实验，检验整体性能、证据保留、记忆复用、预算鲁棒性和故障诊断能力。
\end{enumerate}

\subsection{研究意义}

\subsubsection{理论意义}

第一，本课题有助于把大语言模型智能体研究从“能力堆叠”推进到“治理建模”。现有许多智能体系统强调规划、工具调用、角色协作或记忆模块，但在长时程任务中，真正限制系统稳定性的往往是上下文选择、记忆提升、冲突抑制和预算分配之间的耦合。本课题将这些问题统一为治理策略问题，有助于为长时程智能体系统提供更清晰的问题定义。

第二，本课题有助于区分运行内上下文治理与跨运行长期记忆治理。当前研究中，短期上下文、长期记忆、RAG 检索结果和多智能体共享笔记经常被混合处理，导致机制边界不清。本课题将象群上下文管理策略定位为运行内共享上下文治理，将克拉克星鸦长时记忆图谱定位为跨运行经验沉淀与返回，两者通过 MGCM 治理接口衔接，从而形成更明确的理论分工。

第三，本课题有助于推动智能体长期记忆从平面检索走向关系化记忆。传统 RAG 更多关注文档片段检索，普通长期记忆更多关注历史对话召回，而长时程智能体需要处理任务、证据、时间、修正、冲突和复用范围之间的关系。本课题将长期记忆组织为图谱，有助于提升记忆的可解释性、可审计性和安全返回能力。

第四，本课题有助于建立智能体系统评估中的证据边界意识。当前研究稿件已经明确区分公开 benchmark 主证据与诊断性补充证据，并对 RQ2、RQ3 等未完全闭合的实验结果保持谨慎表述。这种写法有助于形成更符合学术规范的实验叙事，避免把小样本、单 seed 或自定义套件结果过度解释为普遍结论。

\subsubsection{现实意义}

第一，本课题可服务于复杂任务智能体编排框架建设。长篇研究、代码开发、自动实验、数据分析和报告生成等任务都具有跨轮次、跨工具、跨文件和跨子任务的特征。如果没有上下文治理与长期记忆机制，系统容易出现历史约束丢失、重复检索、错误经验复用和任务轨迹不可追踪等问题。本课题的机制可为智能体底层框架提供可落地的设计方案。

第二，本课题可提升大语言模型应用的稳定性与可维护性。实际部署中的智能体系统不仅要能完成任务，还要能解释为什么保留某些上下文、为什么召回某段历史经验、为什么某条记忆被废弃，以及失败发生在哪个步骤。事件中心可观测性和 MGCM 记忆治理能够为调试、审计和复现实验提供结构化证据。

第三，本课题可降低长时程任务的成本和风险。长上下文模型虽然不断发展，但直接堆积全部历史会造成 token 成本、位置偏置、噪声累积和隐私风险。通过预算感知压缩、显著性排序和渐进式长期记忆返回，可以在保持任务连续性的同时减少无效上下文注入，提高系统成本效率。

\section{与选题相关的国内外研究状况及最新动态}

\subsection{大语言模型自主智能体研究}

大语言模型自主智能体通常被建模为在推理、行动、工具调用和环境反馈之间循环的执行系统。ReAct 将显式推理轨迹与工具行动相结合，使模型能够在“思考-行动-观察”的闭环中逐步推进任务；Reflexion 进一步将执行反馈转化为语言式自我改进信号，说明失败经验可以被重新注入后续决策。Wang 等关于大语言模型自主智能体的综述则系统总结了规划、记忆、工具使用、反思和环境交互在智能体系统中的作用。

这一方向为本课题提供了闭环执行基础，但其主要对象仍是单个智能体如何更好地行动、反思和纠错。对于多个智能体在长时间跨度中如何共享有限上下文、如何区分运行内共享记忆和跨运行长期记忆、如何在预算限制下构造提示，现有单智能体闭环研究仍不能直接回答。

\subsection{基于大语言模型的多智能体系统研究}

多智能体系统研究在 LLM 时代迅速发展。传统多智能体理论强调协作、协调、组织结构和通信协议；在大语言模型语境下，研究重点逐渐转向角色画像、通信模式、环境接口、任务分解和能力获得机制。CAMEL 展示了角色扮演智能体协作与数据生成的可行性；AutoGen 将可对话智能体与群聊编排抽象为通用框架；MetaGPT 和 ChatDev 将 SOP、结构化消息和软件开发流水线纳入协作骨架；MASTER、MegaAgent、MACNet 和 Cross-Team Orchestration 进一步探索了专业化角色招募、大规模协作和跨团队编排。

这些研究说明，多智能体系统的性能并不只是由单个模型能力决定，还取决于角色分工、通信协议和协作结构。然而，现有多智能体框架更关注任务分解与团队组织，对“已产生的信息应如何在角色之间传播、哪些上下文应持续保持影响力、哪些经验应提升为未来可召回的长期记忆”等治理问题讨论不足。本课题正是从这一空缺切入。

\subsection{长上下文建模、位置偏置与上下文压缩研究}

长上下文建模为长时程智能体提供了底层条件。围绕 RoPE 扩展、分段位置编码和长窗口训练的研究不断提升模型可接受的 token 长度。但 Lost in the Middle 等研究表明，能够容纳更多 token 并不意味着模型能够稳定利用更多证据，关键信息处于上下文中段时模型性能可能明显下降。这说明长上下文容量与长上下文治理是两个不同问题。

与扩展窗口并行发展的另一条路线是上下文压缩。Selective Context 通过识别并删除冗余内容降低推理成本；LLMLingua 将提示压缩系统化为预算控制、迭代压缩和语义保持问题。这些研究表明，提示构造应是受约束的信息选择，而不是简单截断或顺序堆积。不过，现有压缩方法通常以连续文本为对象，较少显式处理多智能体系统中的异质片段，如工具观测、委派包、共享笔记、长期记忆子图和策略控制片段。本课题拟在此基础上将压缩提升为类型化、阶段感知的运行时治理问题。

\subsection{长期记忆、RAG 与图结构检索研究}

检索增强生成说明外部知识访问能够改善知识密集型任务。随着智能体系统发展，长期记忆逐渐被作为独立于模型参数之外的系统来设计。Generative Agents 将记忆、反思和计划结合起来，展示了记忆机制对群体行为和长期连续性的作用；MemoryBank 引入遗忘曲线和长期记忆更新；MemGPT 借鉴操作系统分层存储思想，将上下文窗口、外部存储和控制流组织为类似虚拟内存的系统；From RAG to Memory 则从非参数持续学习角度讨论了 RAG 向记忆系统演化的方向。

当记忆不再只是平面文本片段，而是包含事件、实体、证据、范围和冲突关系时，图结构检索成为自然选择。GraphRAG 将原始文本组织为实体图和社区摘要，证明图化表示能够支持超越普通 chunk-level RAG 的全局问答与主题级总结。现有 GraphRAG 主要面向静态语料库知识检索，而本课题关注的是智能体系统自身跨运行经验的结构化组织，因此图谱中的节点不是普通文档片段，而是经过治理的记忆胶囊。

\subsection{系统级评测、社会交互与可观测性研究}

智能体评测正在从静态题目分数转向交互行为、长程决策、社会智能和鲁棒性。AgentBench 将评测对象从单步回答扩展到多环境交互；ChatEval 通过多智能体辩论式评审模拟多人裁决；SOTOPIA 和 SOTOPIA-$\pi$ 关注语言智能体在社会情境中的目标完成能力和交互式学习；PEAR 等 planner-executor-agent benchmark 则开始评估规划、执行、评估与修正的系统闭环。

这类研究提示，长时程智能体不能只看最终分数，还需要观察轨迹是否可重建、失败是否可定位、记忆是否正确复用、预算收缩时是否退化。本课题的事件中心可观测性正是针对这一需求设计：每个运行事件记录 run ID、thread ID、task ID、阶段和元数据，使 todo 投影、记忆治理和失败诊断建立在同一套结构化证据之上。

\subsection{国内外研究最新动态与文献述评}

从最新动态看，2024 年以来大语言模型多智能体研究进入快速系统化阶段。一方面，AutoGen、MetaGPT、ChatDev 等系统推动了可对话智能体、软件工程智能体和角色化协作框架的发展；另一方面，AgentBench、SOTOPIA、PEAR 等评测工作推动研究从静态 NLP 指标转向交互式系统能力。2025--2026 年的 MASTER、MegaAgent、MACNet、Cross-Team Orchestration、RAG for AIGC 和 public opinion simulation 等工作继续拓展多智能体协作、规模扩展、检索增强和社会模拟应用。

综合来看，已有研究仍存在四点不足：第一，长时程多智能体系统中的上下文治理与长期记忆治理经常混用，缺少清晰分层；第二，许多长期记忆方法仍以平面文本或向量检索为主，对经验之间的时间、证据、冲突和修正关系建模不足；第三，长上下文和提示压缩研究较少面向多智能体运行时中的异质上下文片段；第四，系统评测中主 benchmark、诊断套件、模块级实验和工程时延实验的证据等级需要更严格区分。因此，本课题拟在现有研究基础上，提出一个边界更清晰、证据更可追踪的上下文治理与长期记忆图谱研究框架。

\section{本课题的基本内容}

\subsection{研究对象与问题定义}

本课题研究对象为长时程多智能体大语言模型系统。系统接收用户请求 $x$，运行在线程状态 $\mathcal{H}_t$ 中，拥有工具集合 $\mathcal{T}$、检索知识底座 $\mathcal{R}$、智能体集合 $\mathcal{A}$ 以及多层记忆存储。在决策步 $t$，系统需要构造可执行提示上下文 $\mathcal{C}_t$，并从推理、工具调用、检索、委派和记忆治理等动作空间中选择动作。

压缩前的候选上下文可以表示为：
\[
\mathcal{X}_t=\mathcal{H}_t \cup \mathcal{M}^{local}_t \cup \mathcal{M}^{collective}_t \cup \mathcal{M}^{nutcracker}_t \cup \mathcal{R}_t \cup \mathcal{D}_t \cup \mathcal{S}_t.
\]
其中，$\mathcal{M}^{local}_t$ 为局部工作记忆，$\mathcal{M}^{collective}_t$ 为受治理的共享记忆，$\mathcal{M}^{nutcracker}_t$ 为长期记忆图谱返回的候选记忆，$\mathcal{R}_t$ 为检索证据，$\mathcal{D}_t$ 为委派数据包和协调元数据，$\mathcal{S}_t$ 为策略指令或结构化控制片段。

提示构造受有限预算 $B_t$ 约束。令每个候选片段 $z_i$ 的长度为 $\ell_i$，则系统必须选择子集 $\mathcal{C}_t \subseteq \mathcal{X}_t$，满足：
\[
\sum_{z_i \in \mathcal{C}_t} \ell_i \leq B_t.
\]
因此，本课题将上下文构造定义为受预算约束的信息选择问题，而不是对话历史的线性追加问题。

\subsection{象群上下文管理策略}

象群上下文管理策略对应论文第三章的核心机制。它并非神经网络内部注意力层，而是一种系统级显著性策略，用于在多智能体协作过程中判断哪些候选上下文更应保留、压缩或传播。对于候选片段 $z_i$，策略分数定义为：
\[
s_i^{EA}=\alpha Rel(z_i,x)+\beta Val(z_i)+\gamma Reuse(z_i)+\delta Out(z_i)-\lambda Red(z_i).
\]
其中，$Rel$ 表示任务相关性，$Val$ 表示验证支持程度，$Reuse$ 表示跨运行或跨智能体复用程度，$Out$ 表示与成功结果的关联，$Red$ 表示冗余惩罚。该策略的核心偏置是：已经验证、可复用且与成功结果有关的上下文，应当优先于只是最近出现但未验证的信息。

该机制主要解决三个问题：第一，在多智能体协作时减少新近但低价值信息对提示空间的占用；第二，在工具观测和委派结果之间保持关键证据连续性；第三，为上下文压缩和跨角色路由提供可解释的排序依据。

\subsection{克拉克星鸦长时记忆图谱机制}

克拉克星鸦长时记忆图谱机制对应论文第四章的核心机制。该机制关注哪些执行经验值得被提升为长期记忆、如何表示为记忆胶囊节点、如何建立关系边，以及未来如何按任务线索召回局部子图。

对于记忆候选 $m$，提升分数定义为：
\[
\mu(m)=\rho_1 Val(m)+\rho_2 Reuse(m)+\rho_3 Out(m)+\rho_4 Gen(m)-\rho_5 Noise(m).
\]
当 $\mu(m)\geq \theta_{promote}$ 时，记忆被提升为长期记忆胶囊并写入图谱。未来任务中，返回分数定义为：
\[
r(m,x,\mathcal{H}_t)=\kappa_1 Match(m,x)+\kappa_2 Thread(m,\mathcal{H}_t)+\kappa_3 Reuse(m)-\kappa_4 Stale(m).
\]
当 $r(m,x,\mathcal{H}_t)\geq \theta_{return}$ 时，该记忆或其局部子图才会被重新引入。

在工程实现上，长期记忆图谱拟作为独立可插拔模块接入智能体运行时。节点采用 \texttt{:MemoryCapsule} 表示，记录任务目标、摘要、结果、知识范围、标签、实体、结构化 claims、显著性、置信度、状态、复用次数和时间戳；边固定为时间相邻、同任务、证据支持、范围重叠、修正替代和显式冲突六类。检索流程采用向量召回、全文召回、scene-aware 重排、一跳子图扩展、陈旧抑制、冲突消解和渐进式披露。

\subsection{长时程注意力、上下文压缩与事件中心可观测性}

长时程注意力用于维持当前运行内部的轨迹连续性。系统维护近期窗口、历史摘要、过程状态快照和产物汇总层，使长任务在原始对话过长时仍能保持“正在做什么、已经做了什么、哪些约束仍有效、哪些产物值得继续使用”的结构化状态。

上下文压缩将异质候选片段转化为预算约束下的优化问题。每个片段具有类型、长度、象群策略分数和阶段感知权重，系统根据当前执行阶段选择最有价值的上下文进入提示。这样可以避免单纯按时间顺序截断，也可以减少低价值工具输出和重复历史对提示空间的占用。

事件中心可观测性用于记录运行事件并投影为 todo 状态。事件流 $\mathcal{E}=(e_1,e_2,\ldots,e_T)$ 不只是日志，而是记忆治理和故障诊断的证据来源。通过事件中心视角，系统可以把执行轨迹、记忆提升依据、上下文选择结果和失败位置统一记录，为后续实验复现和论文分析提供支撑。

\section{本课题的重点和难点}

\subsection{研究重点}

\begin{enumerate}[label=\arabic*.]
  \item 上下文治理与长期记忆治理的边界划分。论文需要明确象群机制解决运行内上下文竞争问题，星鸦机制解决跨运行经验沉淀和返回问题，避免把所有内容都笼统写成“记忆模块”。
  \item 象群上下文管理策略的可解释建模。需要给出任务相关性、验证支持、复用价值、结果关联和冗余抑制等信号的定义，并说明它们如何服务于预算受限的提示构造。
  \item 长期记忆图谱的数据结构与返回路径。重点在于记忆胶囊节点、关系边、索引、混合召回、一跳扩展、陈旧抑制和冲突消解，而不只是把历史记录存入数据库。
  \item 系统级融合。需要解释 MGCM、长时程注意力、上下文压缩和事件中心可观测性如何与两个核心机制衔接，形成完整的\framework。
  \item 实验协议与结论边界。需要严格区分公开 benchmark 主证据、诊断套件、模块级强基线和工程时延测试，避免把证据等级不同的结果混在一起写。
\end{enumerate}

\subsection{研究难点}

\begin{enumerate}[label=\arabic*.]
  \item 实验覆盖成本较高。长时程多智能体任务依赖真实模型调用，跨模型、跨 seed、跨 benchmark 的完整实验需要较高成本和较长运行时间。
  \item 评价指标容易混杂。整体质量分、记忆 recall、复用率、陈旧记忆注入率、冲突消解成功率和 token 开销反映不同目标，不能简单用单一分数概括全部机制。
  \item RQ3 长期记忆图谱证据尚不完全闭合。当前模块级结果显示图谱机制在关系命中、陈旧抑制和冲突消解方面有优势，但在首屏 recall、relevance 和 usefulness 上仍落后于 flat summary memory 强基线，需要在论文中谨慎表述。
  \item 系数与阈值校准存在风险。象群策略过度偏向历史已验证信息可能压制新证据；星鸦图谱提升阈值过低可能积累错误记忆；上下文压缩若冗余惩罚不当可能丢弃低频关键证据。
  \item 工程实现与学术表述需要统一。现有代码和实验资产中保留了象群注意力、星鸦记忆以及 RQ2、RQ3 实验别名等历史命名，论文主文需统一为象群上下文管理策略和克拉克星鸦长时记忆图谱机制，同时在实验部分说明历史别名。
\end{enumerate}

\section{本课题实验设计及创新之处}

\subsection{实验设计}

本课题实验拟分为系统级 benchmark、机制级对照、模块级强基线和工程时延测试四个层次。

\subsubsection{系统级公开 benchmark 实验}

系统级实验围绕五个研究问题展开：

\begin{longtable}{>{\raggedright\arraybackslash}p{0.14\textwidth}>{\raggedright\arraybackslash}p{0.36\textwidth}>{\raggedright\arraybackslash}p{0.38\textwidth}}
  \caption{研究问题与实验设计}\\
  \toprule
  研究问题 & 主要目标 & 评测数据与对照设置\\
  \midrule
  RQ1 & 检验统一框架在长时程协调、多跳推理和工具辅助任务中的整体收益 & GAIA、HotpotQA、MuSiQue；对照单智能体基础版和去除象群机制的多智能体版本\\
  RQ2 & 检验象群上下文管理策略在高干扰上下文中的证据保留能力 & HotpotQA、MuSiQue 高干扰切片；对照去除象群机制的多智能体版本和单智能体长上下文版本\\
  RQ3 & 检验长期记忆机制在 recall、reuse 和 stale memory 控制方面的表现 & LoCoMo、LongMemEval；对照朴素长期记忆版本和去除星鸦机制版本\\
  RQ4 & 检验长时程注意力与上下文压缩在预算收缩下的鲁棒性 & LongBench、InfiniteBench 切片；对照去除上下文压缩和去除长时程注意力版本\\
  RQ5 & 检验事件中心可观测性对轨迹重建和故障定位的作用 & AgentOrch Failure Suite；对照普通日志版本\\
  \bottomrule
\end{longtable}

当前研究稿件中已经给出部分真实运行结果：RQ1 在 GAIA、HotpotQA、MuSiQue 上覆盖 84 次运行；RQ2 在 HotpotQA 和 MuSiQue 高干扰切片上覆盖 97 次运行；RQ3 在 LoCoMo 和 LongMemEval 上覆盖 162 次运行；RQ4 在 InfiniteBench 和 LongBench 切片上覆盖 102 次运行。正式论文阶段将继续补充模型、seed 和样本覆盖。

\subsubsection{长期记忆图谱模块级强基线实验}

为避免系统级结果掩盖长期记忆图谱本身的贡献，本课题设置独立模块级 benchmark。比较对象包括：

\begin{enumerate}[label=\arabic*.]
  \item \texttt{no\_long\_term\_memory}：不返回长期记忆；
  \item \texttt{vector\_memory}：仅使用向量索引返回候选胶囊；
  \item \texttt{flat\_summary\_memory}：将同任务族高状态节点压缩为平面摘要；
  \item \texttt{naive\_graph\_memory}：允许向量/全文混合召回和一跳扩展，但不使用 scene-aware 重排、陈旧抑制和冲突裁决；
  \item \texttt{clarks\_nutcracker\_graph}：完整启用混合召回、关系感知扩展、陈旧抑制、冲突消解和渐进式披露。
\end{enumerate}

评价指标包括 capsule recall@k、subgraph relevance、relation hit rate、returned memory usefulness、stale node injection rate、conflict resolution success 和 latency。当前已有结果显示，完整图谱机制在五个 seed 下 relation hit 均值达到 0.0352，stale rate 为 0，conflict success 均值为 0.9833；但 flat summary memory 在 recall、relevance 和 usefulness 上仍更强。因此，论文将把当前图谱优势定位为结构化关系返回与错误记忆抑制，而不是全面优于所有强基线。

\subsubsection{工程时延测试}

工程时延测试用于验证长期记忆图谱模块是否具有可插拔部署可行性。当前 Neo4j 图规模测试包含 scale-1、scale-2 和 scale-4 三档，节点数量分别为 360、720 和 1440，结果显示 \texttt{recall()} 平均耗时约 91.60--99.08 ms，P95 为 120.24--142.44 ms，\texttt{fetch\_capsule\_details()} 平均耗时约 3.67--3.76 ms。该结果说明图谱模块在当前规模下能够保持交互式响应。

\subsection{技术路线}

本课题技术路线如下：

\begin{quote}
文献梳理与问题抽象
$\rightarrow$
长时程多智能体系统形式化建模
$\rightarrow$
象群上下文管理策略设计
$\rightarrow$
克拉克星鸦长时记忆图谱机制设计
$\rightarrow$
MGCM、长时程注意力、上下文压缩和事件中心可观测性融合
$\rightarrow$
系统级 benchmark 与模块级实验
$\rightarrow$
结果讨论、局限分析与论文定稿。
\end{quote}

\subsection{创新之处}

\begin{enumerate}[label=\arabic*.]
  \item 研究视角创新。将长时程多智能体系统重新表述为上下文、记忆、协作和事件证据的统一治理问题，而不是把智能体能力简单拆成规划、工具和记忆模块。
  \item 机制分层创新。明确区分运行内上下文治理与跨运行长期记忆治理，分别提出象群上下文管理策略和克拉克星鸦长时记忆图谱机制，形成两个可独立成文又能系统融合的章节。
  \item 记忆表示创新。将长期记忆表示为记忆胶囊图谱，显式建模时间、任务、证据、范围、修正和冲突关系，使记忆返回从平面召回转向关系化子图返回。
  \item 评估框架创新。将公开 benchmark、诊断套件、模块级强基线和工程时延实验分层呈现，强调证据等级与结论边界的一致性。
  \item 工程落地创新。长期记忆图谱设计为独立 Neo4j 可插拔模块，可通过标准接口接收候选记忆胶囊并返回可注入提示的局部子图，有利于接入现有 AgentOrch 或其他智能体运行时。
\end{enumerate}

\section{可行性论证}

\subsection{理论可行性}

本课题具有明确的理论基础。自主智能体研究提供推理-行动-观察闭环；多智能体系统研究提供角色协作、通信与组织结构基础；长上下文和提示压缩研究提供预算受限提示构造依据；RAG、长期记忆和 GraphRAG 研究提供外部知识访问和图结构检索基础；系统级评测与可观测性研究提供轨迹诊断和实验设计依据。这些理论能够支撑本课题从问题定义、机制设计到实验验证的完整链条。

\subsection{数据与实验可行性}

当前研究稿件已经包含较完整的实验资产和结果文件，主表、主图和样本覆盖信息可追溯到 \texttt{archive/paper\_benchmark\_suite/report\_assets} 以及正式运行记录。现有公开 benchmark 切片覆盖 GAIA、HotpotQA、MuSiQue、LoCoMo、LongMemEval、LongBench 和 InfiniteBench；诊断套件覆盖 AgentOrch MultiTurn Suite 和 Failure Suite；长期记忆图谱已有模块级强基线结果和 Neo4j 时延测试。因此，论文写作不是从零开始，而是在已有真实结果基础上继续补齐实验覆盖和结论边界。

\subsection{工程可行性}

本课题依托已有 AgentOrch/AgentTorch 风格智能体编排框架与本地实验环境。源稿件中长期记忆图谱模块已经以 Neo4j 5.x 为后端完成可插拔设计，支持 \texttt{MemoryCapsule} 节点、关系边、全文索引、向量索引、混合召回、一跳扩展和详情查询。实验资产中已有正式 benchmark 表格和图像，说明工程管线已经具备基础可复现能力。后续主要任务是继续补充实验覆盖、统一命名、强化统计分析和完善论文叙事。

\subsection{风险控制}

针对研究中的不确定性，拟采取以下控制措施：第一，所有实验结论均标明数据来源、样本数、模型数、seed 数和运行数，避免过度解释；第二，对 RQ2、RQ3 等尚未完全闭合的模块收益保持谨慎表述；第三，将长期记忆图谱模块级结果与系统级 benchmark 结果分开讨论；第四，继续增加强基线和消融实验，尤其关注 flat summary memory 这一强基线；第五，保留事件日志和运行元数据，保证实验失败和异常结果可以追踪。

\section{论文提纲}

拟定论文题目为：\textbf{\reporttitle}。

\subsection{摘要}

摘要部分将概括长时程多智能体大语言模型系统面临的上下文预算、记忆治理、协作连续性和可观测性挑战，说明本文提出的\framework 及其两个核心机制，并总结系统级 benchmark、模块级强基线和工程时延实验的主要发现。

\subsection{目录}

论文目录拟按照以下结构展开：

\subsubsection*{第一章 绪论}
\begin{enumerate}[label=1.\arabic*]
  \item 背景及意义
  \begin{enumerate}[label=1.1.\arabic*]
    \item 研究背景
    \item 研究意义
  \end{enumerate}
  \item 研究现状
  \begin{enumerate}[label=1.2.\arabic*]
    \item 大语言模型自主智能体研究
    \item 基于大语言模型的多智能体系统研究
    \item 长上下文建模与上下文压缩研究
    \item 长期记忆、RAG 与图结构检索研究
    \item 系统级评测与可观测性研究
    \item 文献述评
  \end{enumerate}
  \item 研究方法与内容
  \begin{enumerate}[label=1.3.\arabic*]
    \item 研究方法
    \item 研究内容
    \item 技术路线
  \end{enumerate}
  \item 研究创新与不足
  \begin{enumerate}[label=1.4.\arabic*]
    \item 研究创新
    \item 研究不足
  \end{enumerate}
\end{enumerate}

\subsubsection*{第二章 理论基础与问题建模}
\begin{enumerate}[label=2.\arabic*]
  \item 大语言模型智能体与多智能体系统基础
  \item 长时程任务中的上下文预算约束
  \item 多层记忆系统与记忆治理理论
  \item 长时程多智能体系统的问题形式化
  \item 本章小结
\end{enumerate}

\subsubsection*{第三章 象群启发的多智能体上下文治理方法}
\begin{enumerate}[label=3.\arabic*]
  \item 象群上下文管理策略的设计动机
  \item 候选上下文类型与显著性信号定义
  \item 基于验证、复用和结果关联的上下文排序模型
  \item 预算感知上下文压缩与跨角色委派数据包
  \item RQ2 高干扰上下文实验与消融分析
  \item 本章小结
\end{enumerate}

\subsubsection*{第四章 克拉克星鸦长时记忆图谱机制}
\begin{enumerate}[label=4.\arabic*]
  \item 克拉克星鸦长期记忆机制的设计动机
  \item 记忆候选抽取、提升判定与胶囊化表示
  \item 记忆图谱节点、关系边与索引设计
  \item 混合召回、子图扩展、陈旧抑制与冲突消解
  \item Neo4j 可插拔实现与工程时延测试
  \item 模块级强基线实验与结果分析
  \item 本章小结
\end{enumerate}

\subsubsection*{第五章 系统集成与综合实验验证}
\begin{enumerate}[label=5.\arabic*]
  \item \framework 的整体运行流程
  \item MGCM 记忆治理与事件中心可观测性
  \item 长时程注意力与上下文压缩的系统融合
  \item 实验设置、公开 benchmark 与基线设计
  \item RQ1 长时程任务整体结果
  \item RQ3 长期记忆系统级结果
  \item RQ4 预算鲁棒性结果
  \item RQ5 可观测性与多轮连续性补充分析
  \item 跨模型一致性、统计检验与结果讨论
  \item 本章小结
\end{enumerate}

\subsubsection*{第六章 结论与展望}
\begin{enumerate}[label=6.\arabic*]
  \item 研究结论
  \item 理论贡献与工程启示
  \item 研究局限
  \item 后续研究展望
\end{enumerate}

\subsubsection*{参考文献}

参考文献部分将按学校规定格式统一整理自主智能体、多智能体系统、长上下文、上下文压缩、长期记忆、GraphRAG、系统评测和相关工程框架文献。

\subsubsection*{致谢}

致谢部分将对导师指导、课题组支持、实验环境支持和相关开源项目进行说明。

\subsubsection*{个人简介}

个人简介部分将列出作者教育经历、研究方向、项目经历和阶段性成果。

\section{调研进度安排}

\begin{longtable}{>{\raggedright\arraybackslash}p{0.18\textwidth}>{\raggedright\arraybackslash}p{0.70\textwidth}}
  \caption{课题调研与论文写作进度安排}\\
  \toprule
  时间安排 & 主要任务\\
  \midrule
  2026.04至05 & 完成开题报告定稿，明确论文题目、章节结构、两大核心机制和实验边界；整理现有 \texttt{cas-sc-template-zh.tex}、实验表格和参考文献。\\
  2026.06至07 & 完成第一章和第二章初稿；系统补充自主智能体、多智能体系统、上下文压缩、长期记忆和图结构检索文献。\\
  2026.08至09 & 完成第三章象群上下文治理方法写作；整理 RQ2 高干扰上下文实验，补充消融、指标定义和失败案例分析。\\
  2026.10至11 & 完成第四章克拉克星鸦长时记忆图谱机制写作；完善 Neo4j 图谱实现、模块级强基线实验和工程时延测试。\\
  2026.12至2027.01 & 完成第五章系统集成与综合实验验证；补齐 RQ1--RQ5 的结果分析、跨模型一致性和统计检验。\\
  2027.02至03 & 完成第六章结论与展望；统一全文术语、图表、参考文献、实验边界和负面结果表述。\\
  2027.04至05 & 完成论文查重、格式规范、预答辩修改、正式答辩材料和可能的投稿版论文整理。\\
  \bottomrule
\end{longtable}

\section{预期目标}

本课题预期实现以下目标：

\begin{enumerate}[label=\arabic*.]
  \item 完成长时程多智能体大语言模型系统的问题形式化，将上下文构造、记忆治理、任务委派和事件可观测性纳入统一建模。
  \item 形成象群启发的多智能体上下文治理方法，明确显著性排序、预算压缩和跨角色上下文传递机制。
  \item 形成克拉克星鸦长时记忆图谱机制，完成记忆胶囊表示、图谱关系设计、混合召回、陈旧抑制和冲突消解流程。
  \item 完成 MGCM、长时程注意力、上下文压缩和事件中心可观测性的系统级融合，形成完整的\framework。
  \item 完成系统级公开 benchmark 切片实验、模块级强基线实验和 Neo4j 工程时延实验，形成可追溯的实验结果。
  \item 完成一篇结构完整、边界清晰、实验可复现的硕士学位论文，并将第三章和第四章进一步整理为可独立投稿的小论文雏形。
\end{enumerate}

\section{与本课题相关的主要参考文献}

\begingroup
\small
\setlength{\tabcolsep}{3pt}
\begin{longtable}{>{\centering\arraybackslash}p{0.05\textwidth}>{\raggedright\arraybackslash}p{0.21\textwidth}>{\raggedright\arraybackslash}p{0.44\textwidth}>{\raggedright\arraybackslash}p{0.20\textwidth}}
  \caption{主要参考文献}\\
  \toprule
  序号 & 作者 & 论文名称或文献名称 & 期刊号与出版年月\\
  \midrule
  1 & Vaswani, A. et al. & Attention Is All You Need & NeurIPS, 2017\\
  2 & Devlin, J. et al. & BERT: Pre-training of Deep Bidirectional Transformers for Language Understanding & NAACL-HLT, 2019\\
  3 & Brown, T. B. et al. & Language Models are Few-Shot Learners & NeurIPS, 2020\\
  4 & Ouyang, L. et al. & Training Language Models to Follow Instructions with Human Feedback & NeurIPS, 2022\\
  5 & OpenAI & GPT-4 Technical Report & OpenAI Research, 2023\\
  6 & Touvron, H. et al. & LLaMA: Open and Efficient Foundation Language Models & Meta AI Technical Report, 2023\\
  7 & Touvron, H. et al. & Llama 2: Open Foundation and Fine-Tuned Chat Models & Meta AI Technical Report, 2023\\
  8 & Yao, S. et al. & ReAct: Synergizing Reasoning and Acting in Language Models & ICLR, 2023\\
  9 & Shinn, N. et al. & Reflexion: Language Agents with Verbal Reinforcement Learning & NeurIPS, 2023\\
  10 & Wang, L. et al. & A Survey on Large Language Model Based Autonomous Agents & Frontiers of Computer Science, 2024\\
  11 & Wooldridge, M. & An Introduction to MultiAgent Systems & Wiley, 2009\\
  12 & Guo, T. et al. & Large Language Model Based Multi-Agents: A Survey of Progress and Challenges & IJCAI, 2024\\
  13 & Li, X. et al. & A Survey on LLM-Based Multi-Agent Systems: Workflow, Infrastructure, and Challenges & Vicinagearth, 2024\\
  14 & Li, G. et al. & CAMEL: Communicative Agents for Mind Exploration of Large Scale Language Model Society & NeurIPS, 2023\\
  15 & Wu, Q. et al. & AutoGen: Enabling Next-Gen LLM Applications via Multi-Agent Conversation & COLM, 2024\\
  16 & Hong, S. et al. & MetaGPT: Meta-Programming for a Multi-Agent Collaborative Framework & ICLR, 2024\\
  17 & Qian, C. et al. & ChatDev: Communicative Agents for Software Development & ACL, 2024\\
  18 & Gan, B. et al. & MASTER: A Multi-Agent System with LLM Specialized MCTS & NAACL-HLT, 2025\\
  19 & Wang, J. et al. & MegaAgent: Scaling Large Language Model-Based Multi-Agent Systems without Predefined Workflows & ACL Findings, 2025\\
  20 & Qian, C. et al. & Scaling Large Language Model-Based Multi-Agent Collaboration & ICLR, 2025\\
  21 & Du, Z. et al. & Multi-Agent Collaboration via Cross-Team Orchestration & ACL Findings, 2025\\
  22 & Liu, N. F. et al. & Lost in the Middle: How Language Models Use Long Contexts & TACL, 2024\\
  23 & Li, R. et al. & Extending Context Window in Large Language Models with Segmented Base Adjustment for Rotary Position Embeddings & Applied Sciences, 2024\\
  24 & Li, Y. et al. & Compressing Context to Enhance Inference Efficiency of Large Language Models & EMNLP, 2023\\
  25 & Jiang, H. et al. & LLMLingua: Compressing Prompts for Accelerated Inference of Large Language Models & EMNLP, 2023\\
  26 & Lewis, P. et al. & Retrieval-Augmented Generation for Knowledge-Intensive NLP Tasks & NeurIPS, 2020\\
  27 & Park, J. S. et al. & Generative Agents: Interactive Simulacra of Human Behavior & UIST, 2023\\
  28 & Zhong, W. et al. & MemoryBank: Enhancing Large Language Models with Long-Term Memory & AAAI, 2024\\
  29 & Packer, C. et al. & MemGPT: Towards LLMs as Operating Systems & Technical Report, 2023\\
  30 & Gutierrez, B. J. et al. & From RAG to Memory: Non-Parametric Continual Learning for Large Language Models & ICML, 2025\\
  31 & Edge, D. et al. & From Local to Global: A Graph RAG Approach to Query-Focused Summarization & Microsoft Research, 2024\\
  32 & Brea, J. et al. & Computational Models of Episodic-Like Memory in Food-Caching Birds & Nature Communications, 2023\\
  33 & Liu, X. et al. & AgentBench: Evaluating LLMs as Agents & ICLR, 2024\\
  34 & Chan, C.-M. et al. & ChatEval: Towards Better LLM-Based Evaluators through Multi-Agent Debate & ICLR, 2024\\
  35 & Zhou, X. et al. & SOTOPIA: Interactive Evaluation for Social Intelligence in Language Agents & ICLR, 2024\\
  36 & Wang, R. et al. & SOTOPIA-$\pi$: Interactive Learning of Socially Intelligent Language Agents & ACL, 2024\\
  37 & Dong, Z. et al. & PEAR: Plan, Execute, Assess, and Refine: A Robust Planner-Executor-Agent Benchmark & EACL Findings, 2026\\
  38 & Zhao, P. et al. & Retrieval-Augmented Generation for AI-Generated Content: A Survey & Data Science and Engineering, 2026\\
  39 & Jimenez-Romero, C. et al. & Multi-Agent Systems Powered by Large Language Models: Applications in Swarm Intelligence & Frontiers in Artificial Intelligence, 2025\\
  40 & Lan, H. et al. & Public Opinion Dissemination Simulation Based on Large Language Model Multi-Agent Systems & Scientific Reports, 2026\\
  \bottomrule
\end{longtable}
\endgroup

\section*{说明}

本报告由当前研究源文件 \texttt{cas-sc-template-zh.tex} 转写为硕士论文开题报告。正式提交前，应根据学院模板补充学号、专业、导师、学院等封面信息，并按学校规定将参考文献统一转换为 GB/T 7714 格式。

\end{document}
